\documentclass[11pt,a4paper]{article}
\usepackage[T1]{fontenc}
\usepackage[utf8]{inputenc}
\usepackage[english]{babel}
\usepackage{amsmath,amssymb}
\usepackage[margin=2cm]{geometry}
\setlength{\parindent}{0pt}
\setlength{\emergencystretch}{3em}
\allowdisplaybreaks

\title{About \texttt{partial\_derivatives\_beyond.tex}}
\author{}
\date{}
\begin{document}
\maketitle
This note explains what the test data in \texttt{partial\_derivatives\_beyond.tex} is used for. The file is read by \texttt{test\_multivariable\_calculus.c}.

\begin{enumerate}
\item \textbf{Purpose.} The file holds 6 true statements that the engine does not prove yet. The test checks only that none is called false: the engine may answer that it does not know, never that a true statement is false. When one becomes provable it moves to \texttt{partial\_derivatives.tex}.

\item \textbf{Why each is beyond reach.} Items 1--3 need the exact values of sine and cosine at $\frac{\pi}{3}$, $\frac{\pi}{4}$ and $\frac{\pi}{6}$: the derivatives are computed and confirmed with numbers, but the exact algebra of chapter 11 does not reduce $\cos\frac{\pi}{3}$ to $\frac{1}{2}$ (the same happens to $\cos\frac{\pi}{3} = \frac{1}{2}$ alone). Items 4--6 are derivatives at a point at which the function is not regular (the argument of an absolute value, or a radicand, is zero): the derivative exists there, but the rules of differentiation give no formula, and it should be computed as the limit of the difference quotient.

\item \textbf{Expected result.} No statement is called false; today all six are answered \texttt{NON DIMOSTRATO}.

\item \textbf{How to run it.} From the top folder of the repository:
\begin{verbatim}
cc -std=c11 -O2 -Wall -Wextra -o /tmp/test_multivariable_calculus \
    tests/multivariable_calculus/test_multivariable_calculus.c && \
    /tmp/test_multivariable_calculus
\end{verbatim}
\end{enumerate}
\end{document}
